%-------------------------
% Resume in Latex
% Based off of Jake's Resume Template
% 中文版：请使用 XeLaTeX 编译
%------------------------

\documentclass[letterpaper,11pt]{article}

\usepackage[UTF8]{ctex}
\usepackage{latexsym}
\usepackage[empty]{fullpage}
\usepackage{titlesec}
\usepackage{marvosym}
\usepackage[usenames,dvipsnames]{color}
\usepackage{verbatim}
\usepackage{enumitem}
\usepackage[hidelinks]{hyperref}
\usepackage{fancyhdr}
\usepackage{tabularx}

%----------FONT OPTIONS----------
% 保持模板风格，不额外指定复杂字体
% 如果 Overleaf 中文字体报错，可改用：
% \usepackage[UTF8,fontset=none]{ctex}
% \setCJKmainfont{Noto Serif CJK SC}
% \setCJKsansfont{Noto Sans CJK SC}
% \setCJKmonofont{Noto Sans Mono CJK SC}

\pagestyle{fancy}
\fancyhf{} % clear all header and footer fields
\fancyfoot{}
\renewcommand{\headrulewidth}{0pt}
\renewcommand{\footrulewidth}{0pt}

% Adjust margins - 保持 Jake's Resume 原版设置
\addtolength{\oddsidemargin}{-0.5in}
\addtolength{\evensidemargin}{-0.5in}
\addtolength{\textwidth}{1in}
\addtolength{\topmargin}{-.5in}
\addtolength{\textheight}{1.0in}

\urlstyle{same}

\raggedbottom
\raggedright
\setlength{\tabcolsep}{0in}

% Sections formatting - 保持 Jake's Resume 原版风格
\titleformat{\section}{
  \vspace{-4pt}\raggedright\large\bfseries
}{}{0em}{}[\color{black}\titlerule \vspace{-5pt}]

%-------------------------
% Custom commands - 基本沿用 Jake's Resume 原命令
\newcommand{\resumeItem}[1]{
  \item\small{
    {#1 \vspace{-2pt}}
  }
}

\newcommand{\resumeSubheading}[4]{
  \vspace{-2pt}\item
    \begin{tabular*}{0.97\textwidth}[t]{l@{\extracolsep{\fill}}r}
      \textbf{#1} & #2 \\
      \textit{\small#3} & \textit{\small #4} \\
    \end{tabular*}\vspace{-7pt}
}

\newcommand{\resumeSubSubheading}[2]{
    \item
    \begin{tabular*}{0.97\textwidth}{l@{\extracolsep{\fill}}r}
      \textit{\small#1} & \textit{\small #2} \\
    \end{tabular*}\vspace{-7pt}
}

\newcommand{\resumeProjectHeading}[2]{
    \item
    \begin{tabular*}{0.97\textwidth}{l@{\extracolsep{\fill}}r}
      \small#1 & #2 \\
    \end{tabular*}\vspace{-7pt}
}

\newcommand{\resumeSubItem}[1]{\resumeItem{#1}\vspace{-4pt}}

\renewcommand\labelitemii{$\vcenter{\hbox{\tiny$\bullet$}}$}

\newcommand{\resumeSubHeadingListStart}{\begin{itemize}[leftmargin=0.15in, label={}]}
\newcommand{\resumeSubHeadingListEnd}{\end{itemize}}
\newcommand{\resumeItemListStart}{\begin{itemize}}
\newcommand{\resumeItemListEnd}{\end{itemize}\vspace{-5pt}}

%-------------------------------------------
%%%%%%  RESUME STARTS HERE  %%%%%%%%%%%%%%%%%%%%%%%%%%%%

\begin{document}

%----------HEADING----------
\begin{center}
    \textbf{\Huge 姚泽林} \\ \vspace{8pt}
    \small 华中科技大学计算机技术硕士在读，2028 届 $|$ 求职意向：RAG / 大模型应用研发实习生 \\ \vspace{1pt}
    \small 可立即到岗 $|$
    156-2317-1668 $|$
    \href{mailto:izyyzi@foxmail.com}{\underline{izyyzi@foxmail.com}} $|$
    \href{https://github.com/Gianna-Lynn}{\underline{github.com/Gianna-Lynn}}
\end{center}


%-----------EDUCATION-----------
\section{教育经历}
  \resumeSubHeadingListStart
    \resumeSubheading
      {华中科技大学}{武汉}
      {计算机技术，硕士在读}{2025.09 -- 2028.06 预计毕业}
    \resumeSubheading
      {华中科技大学}{武汉}
      {计算机科学与技术，本科}{2020.09 -- 2025.06}
  \resumeSubHeadingListEnd


%-----------PROGRAMMING SKILLS-----------
\section{技术栈}
 \begin{itemize}[leftmargin=0.15in, label={}]
    \small{\item{
     \textbf{RAG 与向量检索}{：自研 Rust HNSW，SIFT1M benchmark，Qwen3-Embedding，BM25，hybrid retrieval，bge-reranker，answer with citations} \\[2pt]
     \textbf{评测与分析}{：Hit@k，Recall@k，MRR，latency，citation precision / recall，badcase taxonomy，workflow report} \\[2pt]
     \textbf{编程语言}{：Python，Rust，C/C++ 基础} \\[2pt]
     \textbf{AI 基础}{：NumPy，PyTorch 基础，自动微分，计算图，向量表示，Transformer / LLM 应用基础} \\[2pt]
     \textbf{工程工具}{：Git，Linux 基础，Shell 基础，HTTP JSON API，Markdown，JSONL，项目说明文档编写}
    }}
 \end{itemize}


%-----------PROJECTS-----------
\section{项目经历}
  \resumeSubHeadingListStart
    \resumeProjectHeading
        {\textbf{Azusa：自研向量检索核与端到端 RAG 评测系统} $|$ \emph{Rust，Python，HNSW，RAG}}{2026}
        \resumeItemListStart
          \resumeItem{自研 Rust HNSW 向量检索服务，支持向量写入、top-k 查询、metadata 回查与 HTTP JSON 调用；在 SIFT1M 100 万 base 上完成参数 sweep，m8\_efc100、ef\_search=300 下 recall@10=99.40\%。}
          \resumeItem{围绕 HNSW 搜索宽度与建图参数记录 build time、exact@10、latency 等指标；100k 调参落点在 ef\_search=300 下达到 recall@5=97.0\%、avg search time 约 522us。}
          \resumeItem{搭建 Python RAG 链路，覆盖 chunk、embedding、HNSW dense、BM25 sparse、dense-first hybrid 候选融合、rerank、answer with citations 与结果落盘。}
          \resumeItem{实现 brute、hnsw、bm25、hybrid、hybrid\_rerank 五种检索模式；15 条 Stage 2 小规模回归集上，hybrid\_rerank 的 Hit@k、Recall@k、MRR 均为 1.0000。}
          \resumeItem{构建评测 workflow 与 Search Agent v0，记录 query rewrite、多跳 retrieval、evidence fusion、citation 指标，并用 badcase taxonomy 分析 wrong citation / wrong entity 等失败。}
        \resumeItemListEnd

    \resumeProjectHeading
        {\textbf{CMU 10-414/714 Deep Learning Systems 公开课程作业实现} $|$ \emph{Python，NumPy，Autograd}}{2026}
        \resumeItemListStart
          \resumeItem{基于 NumPy 实现简化深度学习框架，覆盖张量数据结构、计算图、自动微分、基础算子与神经网络模块。}
          \resumeItem{实现前向计算、反向传播和参数更新流程，完成线性层、激活函数、损失函数等组件，强化对 PyTorch 底层机制的理解。}
          \resumeItem{结合 CMU 10-414/714 HW0--HW2 进行验证，为后续参与 embedding、reranker、LLM 应用和检索增强系统开发补充底层知识。}
        \resumeItemListEnd
  \resumeSubHeadingListEnd

%-----------AWARDS-----------
\section{荣誉奖项}
  \resumeSubHeadingListStart
    \resumeItem{2025 年全国大学生计算机系统能力大赛~数据库管理系统设计赛~全国总决赛三等奖}
    \resumeItem{第 32 次 CCF-CSP 认证 200 分（前 19\%） $|$ 2024 年美国大学生数学建模竞赛 Honorable Mention $|$ CET-6：549}
  \resumeSubHeadingListEnd

%-------------------------------------------
\end{document}
